
\section{Conclusion}
\label{sec:conclusions}
%recap of results

The present study aimed at reproducing and improving upon the results from \citet{rutter2024}. They showed in their analysis, that twitter sentiment data could be used to empirically validate the CAM and the differences it proposed between anxiety and depression patients regarding their willingness to be exposed to fluctuation in NA levels. According to the CAM, anxiety patients are motivated to maintain a stable NA level in order to not be subjected to strong emotional shifts. This should be also visible in sentiments of published tweets, as measured by variability of NA level estimates.

Using the proposed method of \citet{rutter2024}, the findings could be exactly reproduced using simple mean and standard deviation comparisons between the cohorts. Depression patients showed significantly higher NA levels and NA variability than anxiety patients.

The same conclusions were reached using 6-hour windows, where tweet sentiments within 6-hour windows were aggregated. Here, again, depression patients showed higher NA levels and variability. A key methodological contribution of the present study lies in the extension of prior research through the introduction of windowed aggregation of sentiment data across time. While previous studies in this area have typically relied on aggregated sentiment scores computed across users' entire available time series, our approach applied moving windows to calculate NA levels and NA variability within smaller, temporally localized segments of user activity.

Thus, the previous findings were successfully replicated using these methods, indicating that VADER sentiment scores for tweets can be used to empirically validate the CAM. Showing this on the windowed data indicates that the findings are also robust to short-term fluctuations, that might otherwise negatively influence the validity of results when aggregating data across such extensive time spans.

On all levels of aggregation, the analyzed distributions were shown to not be unimodal symmetric distributions. This is a cause for concern regarding the validity of the results because it means that mean and standard deviation are not the only reasonable measures of NA levels and NA variability respectively. 

Furthermore, a novel approach for validating the CAM was presented. It was shown how LMMs can be used to add an even more robust perspective to this. The LMM has the ability of taking into account many external factors that might influence NA and its fluctuations, that are independent of cohort membership and might otherwise unexpectedly confound the results. Also, it can estimate differences in NA levels through estimating a regression coefficient for one of the cohorts. This also supported the findings from \citet{rutter2024}, showing that depression patients have higher NA levels than anxiety patients.

Another strong addition of the LMM approach is the ability to model NA variability in a different way. While NA variability is operationalized directly as standard deviation of NA scores in the prior approaches, the LMM adds a cleaner way of analyzing variability. We proposed, that NA variability can be operationalized as the amount of variance in NA levels that cannot be explained by external factors or cohort membership as predictors of NA level. Thus, residual variances were analyzed and compared. This approach also managed to support the previous findings, showing that residual variances were higher for the depression cohort, indicating that it is more difficult to model NA levels for depression patients, indicating a stronger proclivity to (seemingly) random fluctuations.

In addition to providing further evidence for the previously published findings, the LMMs indicated strong autocorrelation in tweet sentiments, showing how consecutive tweets often belong to similar periods with regards to affect levels. When investigating all tweets separately, women showed lower NA variability than men, although the effect switched for 6-hour windowed data. Thus, women appeared to have fewer short-term (e.g. single-tweet) mood fluctuations, while they generally have higher NA variability for more smoothed periods.

The inherent coupling between standard deviation and mean was identified as a possible weakness to the methodology, since it could have been that the differences are simply due to the differences in NA level between the cohorts. This characteristic of bounded distributions was analyzed and the findings were found not to have been impacted by it. Even when controlling for mean NA levels, the A cohort members showed higher levels of NA variability.

Furthermore, this study highlights critical challenges in applying sentiment analysis tools to CAM research. While VADER provided stable and interpretable continuous sentiment scores, its lexicon-based approach comes with well-known limitations, including reduced sensitivity to context, sarcasm, and complex language structures prevalent on social media platforms. The robustness of the analyses was checked by replacing the lexicon-based VADER sentiment analysis with a context-aware transformer model (RoBERTa) in order to examine whether the effects hold up for more sophisticated sentiment extraction. In this way, further support was found for the differences in NA levels between the cohorts. All aggregation approaches (tweet-level data, 6-hour windows, daily aggregation) as well as the LMM showed that the depression cohort had higher NA levels than the anxiety cohort. 

Therefore, very strong support was found for the idea that depression patients have generally elevated levels of NA as compared to anxiety patients.

For NA variability, a different pattern emerged. While the simple comparisons of within-user standard deviations did not exhibit any significant differences between the cohorts, the LMM showed opposite results. The residual variance was significantly higher for the anxiety cohort than for the depression cohort, which opposes all previous findings, as well as all other results from this study. This indicates, that the previous results cannot be replicated using RoBERTa for sentiment estimation. 

This is a striking finding, since it suggests, that the previous findings diminish (and possibly even get inverted) when using a more sophisticated sentiment analysis tool which was explicitly made for analyzing Twitter data. 

Thus, using VADER NA scores, the previous findings could confidently be supported. This not only holds for the established methods, but also for the proposed windowing as well as LMM-based analysis. The present study found strong arguments against the status quo while implementing more sophisticated transformer-based sentiment analyses, indicating that the CAM cannot be empirically validated using Twitter data with transformer-based (or at least RoBERTa-based) sentiment estimation. 

Moreover, advancing CAM-based analyses in digital contexts will benefit from continued innovation in the temporal modeling of affect, including the use of windowed aggregation, dynamic modeling approaches, and longitudinal designs. Such methodological developments are essential for capturing the richness and complexity of real-world affective experiences as expressed in social media behavior.
